% Chapter 10: Test-Driven Development and Quality Assurance
% File: chapters/chapter-14-testing/chapter-14-testing.tex

\chapter[Test-Driven Development]{Test-Driven Development and Quality Assurance: Nine Real Drops}
\label{ch:chapter-14}
\chapthumb{10}

\begin{quote}
    \textit{``A test that passes while the thing it tests is broken is worse than no test at
    all: it converts an open defect into a certified one.''}
    \\
    -- the rule this chapter was written to satisfy
\end{quote}

Embedded software is usually tested in two places: on a host, where the compiler is friendly
and the hardware is absent, and on the target, where the opposite is true. Most projects do
only the second, call it ``testing on hardware'', and never learn that the numbers they
trusted were produced by an instrument that was lying.

This chapter follows one feature from first line to validated behaviour: a fall detector for
a wrist-worn or carried device, built on a six-axis IMU. It is a good case study because it
cannot be judged by inspection — a fall is an event in time, and the only way to know whether
you detect it is to drop the thing and look. Along the way, three of the measurements used to
validate it turn out to be wrong in ways that would have certified a defect as fixed.

\minitoc

\clearpage

\section{What Can Be Tested Off Target}
\label{sec:off-target-testing}

The first design decision of a testable embedded feature is not a test at all: it is where
the code lives.

\subsection{Split the Policy From the Peripheral}

A fall detector is two things wearing one name:

\begin{itemize}
    \item \textbf{Acquisition} — power the IMU, configure its scale and output rate, read a
          12-byte burst, convert to milligees. This only exists on the target.
    \item \textbf{Classification} — decide, from a stream of magnitudes and the time between
          them, that what just happened was a fall. This is arithmetic and state, and it does
          not know that a sensor exists.
\end{itemize}

If both live in the driver crate, the second becomes untestable: building the tests drags in
a PAC, a linker script and a target triple, so nobody writes them. This project keeps the
classifier in its own crate with \textbf{no dependencies at all}:

\begin{minted}{toml}
[package]
name = "motion-detection"
version = "0.1.0"
edition = "2021"

[dependencies]
# deliberately empty: no PAC, no HAL, no panic handler
\end{minted}

\begin{tipbox}
    A dependency-free classifier crate is not an aesthetic choice. It is what makes
    \codeinline{cargo test} work on the developer's machine in under a second, which is what
    makes writing the tenth test cheap enough to bother.
\end{tipbox}

\subsection{The API That Makes Testing Possible}

The crate takes samples and returns events, with the sample interval as a parameter rather
than something it looks up:

\begin{minted}{rust}
pub struct FallDetector { /* phase, peak, timers */ }

impl FallDetector {
    pub const fn new() -> Self;
    /// Feed one sample; the interval is microseconds.
    pub fn update_us(&mut self, accel_mg: [i32; 3], dt_us: u32) -> Option<Event>;
    /// True while the device is believed to be in the air.
    pub fn in_free_fall(&self) -> bool;
}

pub enum Event {
    FreeFall { ms: u32 },
    Impact { peak_mg: i32 },
    Fall { peak_mg: i32, freefall_ms: u32 },
}
\end{minted}

Two details are worth copying into any similar design:

\begin{enumerate}
    \item \textbf{Time is an argument.} A detector that calls a clock cannot be fed a
          recording at ten times real speed, and cannot be tested for how it behaves when a
          sample arrives late. Passing the interval makes both possible and costs nothing.
    \item \textbf{Microseconds, not milliseconds.} At 400\,Hz a sample is 2.5\,ms; rounding to
          whole milliseconds is a 20\,\% error in every timeout. This is not a theoretical
          concern — see Section~\ref{sec:instrument-lied}.
\end{enumerate}

\subsection{The Shape of the Whole Test Suite}

Nine host tests cover the classifier. They are not nine variations of the same case; each one
exists because a specific wrong implementation would pass the others:

\begin{table}[H]
\centering
\begin{tabular}{p{6.6cm}p{7.4cm}}
\hline
\textbf{Test} & \textbf{What it would catch} \\
\hline
\codeinline{quiet_board_reports_nothing} & an implementation that reports on any movement \\
\codeinline{a_dip_too_short_is_not_a_fall} & counting a single low sample as free fall \\
\codeinline{impact_without_free_fall_is_reported_as_impact} & calling every knock a fall \\
\codeinline{a_fall_caught_before_landing_is_not_reported} & ignoring the landing entirely \\
\codeinline{handling_is_not_a_fall} & thresholds tuned only to the positive case \\
\codeinline{real_drop_with_impact_after_a_one_g_interlude_is_detected} &
  requiring the impact at the transition \\
\codeinline{real_drop_with_a_long_interlude_is_detected} & the same, with a longer gap \\
\codeinline{a_soft_landing_after_a_short_free_fall_is_detected} &
  a settle rule that is too strict \\
\codeinline{isqrt_is_exact} & square roots that drift and shift every threshold \\
\hline
\end{tabular}
\caption{The classifier's host tests, each justified by the wrong implementation it rejects.}
\label{tab:classifier-tests}
\end{table}

\begin{remark}
    \textbf{Every one of those tests was added after a real failure, not before.} The suite
    started with five and grew as the recordings exposed assumptions: the interlude tests came
    from two drops that were missed, the soft-landing test from a third. That is normal. What
    is not acceptable is deleting a failing test because the new code is ``better''.
\end{remark}

\section{Tests That Pin a Threshold to a Measurement}
\label{sec:tests-pin-thresholds}

A threshold with no test is a guess that compiles. The discipline that changes it is simple:
\textbf{every constant is justified by a number, and every number has a test on the other side
of it.}

\subsection{The Near-Miss Rule}

Any rule that suppresses or classifies something needs a case that comes \emph{close} and
must still behave correctly. For a fall detector the near misses are the interesting ones:

\begin{table}[H]
\centering
\begin{tabular}{llll}
\hline
\textbf{Boundary} & \textbf{Value} & \textbf{Must trigger} & \textbf{Must not trigger} \\
\hline
free fall & $<$ 400\,mg & real drops: 32--50\,mg & shaking: never below 630\,mg \\
free-fall length & $\geq$ 40\,ms & 42, 70, 80, 100, 138, 145, 176--252\,ms & dips of 10--25\,ms \\
landing & $\geq$ 1100\,mg & 1148--12823\,mg & caught in the hand: $\approx$ 990\,mg \\
stillness & $<$ 1500\,mg for 600\,ms & 1.04--1.17\,s observed & bouncing: $>$1500\,mg \\
\hline
\end{tabular}
\caption{Every constant, the measurements on both sides of it, and the control that must not fire.}
\label{tab:thresholds-near-miss}
\end{table}

Two of these margins are uncomfortably thin, and saying so is part of the test:

\begin{warningbox}
    The landing threshold is 1100\,mg with the softest real landing at 1148\,mg — an 8\,\%
    margin, on nine drops, performed by one person on one desk. It is enough to ship a rule
    and not enough to claim the rule is right. The test encodes the measurements, not a
    belief about all landings everywhere.
\end{warningbox}

\subsection{Writing the Test From the Recording, Not From the Idea}

The failing pattern to avoid is a test that restates the implementation:

\begin{minted}{rust}
// Useless: this asserts that the code does what the code does.
#[test]
fn detects_a_fall() {
    let mut det = FallDetector::new();
    assert!(det.update([0, 0, 0], 5).is_some());   // whatever the code says
}
\end{minted}

The useful version uses a sequence of samples taken from a recording, and asserts on the
number a person measured:

\begin{minted}{rust}
#[test]
fn a_soft_landing_after_a_short_free_fall_is_detected() {
    // Drop 3 of the 400 Hz round: 42 ms of free fall, landing of only 1493 mg,
    // then quiet. This is the case an 800 ms settle rule lost.
    let mut det = FallDetector::new();
    feed(&mut det, [990, 40, 60], 300);      // resting
    feed(&mut det, [30, 25, 20], 42);        // in the air
    feed(&mut det, [1200, 700, 400], 20);    // 1493 mg landing
    match feed(&mut det, [980, 40, 60], 700) {
        Some(Event::Fall { peak_mg, .. }) => {
            assert!((1200..1700).contains(&peak_mg), "peak {peak_mg}");
        }
        other => panic!("expected a fall, got {other:?}"),
    }
}
\end{minted}

Three properties make this test worth its lines: the numbers come from hardware, the
assertion range is wider than the measured value (sensor noise is real), and it fails loudly
with the actual event when the behaviour changes.

\subsection{Embed the Recording in the Image}

Host tests prove the logic. They do not prove that the code on the device behaves the same,
because the device may run a different build, a different clock, a different optimisation
level. The cheap bridge between the two is to embed the recordings in the firmware and replay
them at boot:

\begin{minted}{text}
=== self test at boot, no host attached ===
real=clean  synth=ok  200Hz=2/2  400Hz=4/4
\end{minted}

That line is the whole argument: the handling recording produces nothing, the synthetic fall
is detected, and six of the real drops — replayed through the same binary that runs live — are
all found. A regression in the shipped image shows up as a different line, on a device that
needs no laptop.

\begin{tipbox}
    Embed the recordings as \codeinline{const} arrays (a few kilobytes) and replay them
    \emph{before} the main loop. It costs one flash write and turns every power-up into a
    regression test.
\end{tipbox}

\section{Nine Real Drops}
\label{sec:nine-real-drops}

The detector was validated by dropping the board, by hand, and reading what it recorded. Not
because that is rigorous — it is not, and the chapter says so below — but because it is the
only measurement that can distinguish a plausible rule from a working one.

\subsection{The Protocol}

Three rounds, each recorded by the board itself, plus one control:

\begin{enumerate}
    \item Configure the IMU for a range that will not clip ($\pm$8\,g) and the fastest output
          rate it offers (416\,Hz).
    \item Sample as fast as the loop allows, and record every sample \emph{with its interval}.
    \item Drop the board from about 40\,cm, several times, two seconds apart.
    \item Then perform the negative case: shake hard, set it down suddenly, shake again.
    \item Read the log over the debug probe, with a script, and analyse it on the host.
\end{enumerate}

\begin{notebox}
    \textbf{The control is not optional.} A detector that fires on everything detects every
    fall. Every round of drops in this chapter was recorded together with a handling round
    that had to stay silent, and the thresholds were chosen to satisfy both at once.
\end{notebox}

\subsection{What the Recordings Showed}

\begin{table}[H]
\centering
\begin{tabular}{llll}
\hline
\textbf{Round} & \textbf{Rate} & \textbf{Drops} & \textbf{Detected} \\
\hline
1 & 200\,Hz & 2 & 2 \\
2 & 200\,Hz & 5 & 3 (two lost their free fall between samples) \\
3 & 400\,Hz & 4 onto a bed & 4 \\
control & 200\,Hz & handling & 0 false alarms \\
\hline
\end{tabular}
\caption{Three rounds of drops and one control, all recorded on the board.}
\label{tab:drop-rounds}
\end{table}

Round 3 is the one worth reading in full, because it is the round that corrected the design:

\begin{table}[H]
\centering
\begin{tabular}{llll}
\hline
\textbf{Free fall} & \textbf{Landing} & \textbf{Quiet before the verdict} & \textbf{Verdict} \\
\hline
176\,ms & 11699\,mg & 1.04\,s & fall \\
223\,ms & 13840\,mg & 1.05\,s & fall \\
248\,ms & 13149\,mg & 1.14\,s & fall \\
252\,ms & 13742\,mg & 1.17\,s & fall \\
\hline
\end{tabular}
\caption{Four drops onto a bed, measured. The landings are three to four times harder than
the prediction that a mattress would absorb them.}
\label{tab:bed-round}
\end{table}

\subsection{The Three Assumptions the Data Destroyed}

\paragraph{``A fall is a big impact.''} Real landings measured 1148 to 12823\,mg. Handling
peaked at 1801\,mg. The ranges overlap: a board laid down flat lands more gently than a board
shaken hard, so \textbf{no impact threshold can separate a fall from a knock}. What separates
them is the free-fall phase. The impact threshold only has to distinguish ``landed'' from
``returned to one g because it was caught''.

\paragraph{``The impact follows the free fall immediately.''} In the first recordings it came
120\,ms later, and in the second 400\,ms later. A state machine that required the landing at
the transition detected \textbf{zero of two} drops. Opening a 600\,ms window took the same data
to two of two.

\paragraph{``A bed is soft.''} Predicted: a mattress landing would fall below the 1100\,mg
threshold and be missed. Measured: 11.7 to 13.8\,g. The prediction was wrong, the rule was
right, and the honest conclusion is that nine drops on one bed say nothing about mattresses
in general.

\begin{warningbox}
    \textbf{What this validation is not.} Nine drops, one board, one desk, one bed, one
    person. It proves the detector works on this setup and that the thresholds are not
    arbitrary. It does not prove a fall-detection rate for a product, and a claim that it did
    would be the kind of overreach that gets someone hurt.
\end{warningbox}

\subsection{What Is Still Missing}

Two things are absent from this validation and both are visible in the code:

\begin{itemize}
    \item \textbf{The gyroscope is read and ignored.} Rotation would separate ``dropped while
          flat'' from ``tipped over on a table'', and the data is already in the same burst.
    \item \textbf{A person falling is not a board being dropped.} A carried device falls with
          its owner: shorter free fall, cushioned landing, and a body that moves afterwards.
          The classifier is tuned to the harsher case for the impact threshold and the easier
          one for the free-fall phase, which is a defensible starting point and not an answer.
\end{itemize}

\section{The Instrument That Lied}
\label{sec:instrument-lied}

Everything in the previous section was, at one point, wrong in the same way: the measurements
were taken on a time axis that did not exist. This section is the most important one in the
chapter, because the failure was silent, it survived a code review, and it certified a defect
as fixed.

\subsection{The Sampler That Claimed 400\,Hz and Ran at 158}

The loop was written to sample at 400\,Hz and handed the detector a constant:

\begin{minted}{rust}
const SAMPLE_PERIOD_CYCLES: u32 = 47_500;  // 400 Hz at 19 MHz
const DT_US: u32 = 2_500;                  // 400 Hz, in microseconds
...
application.sample_us(motion.accel_mg, DT_US);
\end{minted}

The loop was paced correctly — it never sampled \emph{faster} than 400\,Hz. But each iteration
also did work: a 12-byte I2C burst, a detector update, a record write, and periodically a
1090-byte display flush over the same bit-banged bus. The true interval was set by the
\textbf{slowest} part of the body, not by the period:

\begin{minted}{text}
measured: 317 samples in 2 s  ->  158 Hz, mean interval 6.3 ms
\end{minted}

Every timeout in the detector was therefore 2.5 times too short, and every duration in the
first analysis — ``40 ms of free fall'' — was 2.5 times too small. The measurements were
self-consistent, which is exactly why nobody noticed.

\subsection{What Gave It Away}

Not a test. A pattern in the output:

\begin{minted}{text}
AIRBORNE  free 85ms   ->  FALL   after 0.62 s
AIRBORNE  free 52ms   ->  FALL   after 0.62 s
AIRBORNE  free 120ms  ->  FALL   after 0.62 s
AIRBORNE  free 40ms   ->  FALL   after 0.53 s
\end{minted}

Four different drops, four different free falls, and the settle time is the same to the
millisecond. \textbf{The number was describing the code, not the world.} That is the signature
of a broken instrument: when an output is constant across inputs that should vary, you are
reading your own parameter back.

\subsection{Three Fixes, and One of Them Was Deleted by the Compiler}

\begin{enumerate}
    \item \textbf{Measure the interval.} The loop takes the difference of the DWT cycle
          counter between samples and passes that everywhere — detector, recorder, alarm.
          \begin{minted}{rust}
let dt_us = now.wrapping_sub(last_sample_cycles) / CYCLES_PER_US;
last_sample_cycles = now;
application.sample_us(motion.accel_mg, dt_us);
          \end{minted}
    \item \textbf{Record the interval with every sample.} A dataset whose timing is unknown
          can only be analysed dishonestly. One byte per sample, in 250\,$\mu$s units.
    \item \textbf{Get the display out of the critical path.} Refresh every 800 samples instead
          of 80, and never while the device is airborne or within 1.2\,s of an event: a flush
          in the middle of a fall stalls the sampler exactly when the samples matter. Measured
          effect: 158\,Hz to 235\,Hz without touching the sensor.
\end{enumerate}

And then the second-order failure, which is the reason this section exists:

\begin{warningbox}
    \textbf{The compiler deleted the per-sample timing array.} The recording was written to
    ``carry its own timing'', and nothing in the crate ever read that array — so the optimiser
    removed it from the image entirely. The code compiled, the comment claimed a property, the
    symbol was absent from the binary, and the recording carried nothing. It was found by
    running \codeinline{arm-none-eabi-nm} and looking for the symbol instead of trusting the
    comment. The writes are volatile now, and the check is part of the habit: \emph{a property
    you cannot see in the artefact is a property you do not have.}
\end{warningbox}

\subsection{The Instrument Rules}

\begin{keytakeaways}
    \item \textbf{Never pass the interval you asked for; pass the one you measured.} The
          intended rate is a wish. The DWT counter is the fact.
    \item \textbf{A constant output across varying inputs means you are measuring yourself.}
          The 0.62\,s settle was a parameter echoed back.
    \item \textbf{A recording must carry its timing.} Otherwise every later analysis inherits
          an assumption nobody remembers making.
    \item \textbf{Verify the artefact, not the intent.} Check that the symbol is in the binary,
          that the counter advanced, that the file grew. Comments describe intentions; the
          binary is what runs.
    \item \textbf{Calibrate the instrument before calibrating the product.} The rate, the
          units and the resolution have to be right before any threshold on top of them means
          anything.
\end{keytakeaways}

\section{What To Take Away}
\label{sec:testing-takeaways}

\begin{keytakeaways}
    \item \textbf{Put the policy in a crate with no dependencies.} If the classifier needs a
          PAC to compile, its tests will not be written.
    \item \textbf{Pass time in as an argument.} It makes fast replay, late samples and
          deterministic tests free.
    \item \textbf{Every threshold is a measurement with a test on both sides.} A rule without
          a near-miss case is a guess with a comment.
    \item \textbf{Embed the recordings and replay them at boot.} Every power-up becomes a
          regression test on the exact binary that ships.
    \item \textbf{Calibrate the instrument first.} A wrong sample rate invalidates every
          number downstream, silently and consistently.
    \item \textbf{Say what the validation is not.} Nine drops prove a rule is not arbitrary;
          they do not prove a detection rate for a product.
\end{keytakeaways}

\subsection{Exercises}

\begin{exercises}
    \item \textbf{(Easy)} Take the classifier out of the driver crate in a project of your own
          and make its tests run with \codeinline{cargo test} on the host. Report how long the
          suite takes before and after.
    \item \textbf{(Easy)} Add a test for the boundary on each side of one threshold: a sample
          exactly at the limit, and one a single unit below it. Which side does your
          implementation treat as inclusive, and is that what the documentation says?
    \item \textbf{(Medium)} Instrument your own sampling loop with \codeinline{DWT->CYCCNT} and
          report the real rate and the worst-case interval. Compare with the rate you intended.
    \item \textbf{(Medium)} Add the gyroscope to the classifier: compute the rotation during the
          free fall and use it to separate ``dropped flat'' from ``tipped over''. What
          additional recording do you need to set the threshold, and can you get it without a
          second person?
    \item \textbf{(Hard)} Design the validation for a detector of a different event of your
          choice (a door being forced, a machine starting, a container being opened) following
          the protocol of this chapter: rounds, a control that must stay silent, near-miss
          cases, and an explicit statement of what the sample does not prove.
\end{exercises}

\begin{notebox}
    \textbf{Where this leads.} The next chapter takes the same evidence-driven approach to the
    layer above the drivers: the board support package, where the question stops being ``does
    this device work'' and becomes ``does this board work, and how do I prove that nobody has
    to know which pin it is on''.
\end{notebox}

